% ============================================================
\chapter{Modul 6: Dictionaries}
% ============================================================

\section{Lektion 18: Schlüssel-Wert-Paare}

\begin{analogie}
Ein Wörterbuch: Du suchst ein Wort (Schlüssel), du findest die Bedeutung (Wert). Nicht nach Position -- nach Name. Genau das braucht man für Sensordaten, Konfiguration und Zustandsbeschreibungen.
\end{analogie}

\begin{lstlisting}
# Roboter-Zustandsobjekt
roboter = {
    "name":          "MobileRobot-1",
    "akku":          87,
    "geschwindigkeit": 0.3,
    "position":      {"x": 1.5, "y": 0.8},
    "sensoren_ok":   True,
    "modus":         "patrouille"
}

# Zugriff
print(roboter["name"])
print(roboter["position"]["x"])

# Sicherer Zugriff (kein Fehler wenn Key fehlt)
temp = roboter.get("temperatur", "nicht verfuegbar")
print(f"Temperatur: {temp}")

# Wert aendern
roboter["akku"] -= 5
roboter["position"]["x"] += 0.1

# Neuen Schlüssel hinzufügen
roboter["letzte_aktion"] = "vorwaerts_fahren"
\end{lstlisting}

\begin{merksatz}
\texttt{dict.get(key, default)} ist der sichere Weg: Wenn der Schlüssel nicht existiert, wird \texttt{default} zurückgegeben statt einer Fehlermeldung. In Robotik-Code immer \texttt{.get()} bevorzugen.
\end{merksatz}

\section{Lektion 19: Sensordaten strukturieren}

\begin{lstlisting}
# Mehrere Sensoren als Dictionary
sensordaten = {
    "ultraschall": {"wert": 0.42, "einheit": "m",   "ok": True},
    "ir_links":    {"wert": 0.18, "einheit": "m",   "ok": True},
    "ir_rechts":   {"wert": 0.35, "einheit": "m",   "ok": True},
    "imu_winkel":  {"wert": 45.2, "einheit": "deg", "ok": True},
    "motortemp":   {"wert": 68.0, "einheit": "C",   "ok": True},
}

# Alle defekten Sensoren finden
print("Sensor-Status:")
for name, daten in sensordaten.items():
    status = "OK" if daten["ok"] else "FEHLER"
    print(f"  {name:15}: {daten['wert']:6.2f} {daten['einheit']} [{status}]")

# Konfiguration aus Dictionary laden
konfig = {
    "max_geschwindigkeit":  0.5,
    "min_abstand":          0.20,
    "akku_warnschwelle":    20,
    "scan_frequenz_hz":     10,
}

def fahre_sicher(abstand, konfig):
    if abstand < konfig["min_abstand"]:
        return 0.0
    return konfig["max_geschwindigkeit"]
\end{lstlisting}

\begin{ros}
In ROS2 werden Parameter genau so gespeichert: als Dictionary-ähnliche Struktur. \texttt{self.declare\_parameter('min\_abstand', 0.20)} ist das ROS2-Äquivalent von \texttt{konfig["min\_abstand"] = 0.20}.
\end{ros}

\begin{uebung}[title=Projektaufgabe Modul 6: Roboter-Konfigurationsmanager]
Erstelle ein Dictionary \texttt{roboter\_config} mit mindestens 8 Parametern. Schreibe Funktionen: \texttt{lade\_config(datei\_name)} (simuliert, gibt Dictionary zurück), \texttt{zeige\_config(config)}, \texttt{aktualisiere\_param(config, key, wert)} (mit Validierung: Wert muss im sinnvollen Bereich liegen).
\end{uebung}
